\begin{mistake}{2026-04-09}{23}

\begin{problemcontent}
以下关于函数极限的命题，正确的是：
(A) 若 \( \lim_{x \to x_0} f(x) \) 存在，则 \( f(x) \) 在 \( x_0 \) 的某去心邻域内有定义
(B) 若 \( f(x) \) 在 \( x_0 \) 处连续，则 \( f(x) \) 在 \( x_0 \) 的某邻域内连续
(C) 若 \( x_0 \) 是 \( f(x) \) 的间断点，则 \( \lim_{x \to x_0} f(x) \) 不存在
(D) 若 \( f(x_0) \) 不存在（无定义），则 \( \lim_{x \to x_0} f(x) \) 不存在
\end{problemcontent}

\begin{solutioncontent}
答案为 (A)。

(A) 正确：极限 \( \lim_{x \to x_0} f(x) \) 的定义前提即要求函数必须在 \( x_0 \) 的某去心邻域内有定义。
(B) 错误：单点连续不能推导局部连续。反例：狄利克雷函数变体（有理数处为 \( x \)，无理数处为 \( 0 \)）仅在 \( x=0 \) 处连续。
(C) 错误：若 \( x_0 \) 为可去间断点，其极限依然存在（如 \( \frac{x^2-1}{x-1} \) 当 \( x \to 1 \) 时）。
(D) 错误：极限是否存在只与去心邻域内的函数行为有关，与该点处的函数值存在与否无关（如 \( \frac{\sin x}{x} \) 当 \( x \to 0 \) 时）。
\end{solutioncontent}

\mistakeknowledge{
\begin{itemize}
 \item \textbf{核心定理}：极限的存在性条件与定义前提（去心邻域内必须有定义）。
 \item \textbf{易错点}：主观混淆了“一点处连续”与“该点邻域内连续”的概念差异。
 \item \textbf{关联知识}：可去间断点的本质（极限存在但无定义或不等于函数值）；特殊函数反例的构造（狄利克雷函数系列）。
\end{itemize}}

\end{mistake}
